\section{Robots humanoides}

\subsection{Marcha bípeda y control de cuerpo completo}

El ponente presentó brevemente los avances de RL aplicado a robots humanoides:
\begin{itemize}
    \item Marcha bípeda estable
    \item Coordinación de cuerpo completo (caminar y manipular al mismo tiempo)
    \item Robustez ante perturbaciones de fuerzas externas
\end{itemize}

\begin{knowledgebox}{Retos particulares de los robots humanoides}
En comparación con los robots cuadrúpedos, el entrenamiento con RL de robots humanoides enfrenta retos mayores:
\begin{itemize}
    \item Apoyo bípedo inestable (margen de estabilidad más estrecho)
    \item Más grados de libertad (más de 30 articulaciones en todo el cuerpo)
    \item Las extremidades superiores e inferiores deben coordinarse
    \item Las caídas tienen consecuencias más graves (el hardware es costoso)
\end{itemize}
\end{knowledgebox}

\subsection{Resumen de la sección}

Los robots humanoides representan el reto de frontera de RL for robotics; los avances actuales son rápidos, pero todavía falta para lograr un despliegue estable.
